% ===================================================================
%  TABEL Nr. 8 — De oogst. — De jacht, de wijnoogst, het vissen.
%  Traduction 2026 d'apres texto/io, controlee sur texto/fr et
%  texto/es. Consigne : voir texto/nl/10-tabel-01.tex.
%
%  L'appel de note est dans le TITRE de la scene, non dans un alinea.
%  LE RENVOI (13) DE LA FAUX MANQUE A L'IDO ; la colonne le rend :
%  ecart declare dans renvoji.py.
% ===================================================================

%%K t08-noto-1 noto
\VUnotes{143.2mm}{%
\textit{(*) Omdat dat woord niet nauwkeurig is : gaat het om vlasoogst,
wijnoogst, appeloogst? Misschien zou het goed zijn «} \VUgras{mesono}
\textit{» te gebruiken dat in Mondo XIV, blz. 242 voorgesteld werd. Maar tot nu
toe is die nieuwe wortel niet door de Academie aangenomen en wacht hij op
besprekingen volgens het gebruikelijke idistische reglement.}%
}

%%K t08-apar-1 apar
\VUsaut{5.0mm}
\VUtitre{15.0pt}{60}{TABEL Nr. 8}
\VUsaut{3.0mm}
\VUcentre{13.2pt}{90}{\textit{Eerste tafereel.}}
\VUsaut{3.0mm}
\VUpk{10.2pt}{40}{De oogst (*).}
\VUsaut{4.0mm}
\VUpk{10.2pt}{40}{De aanblikken van het land.}
\VUsaut{3.0mm}

%%K t08-c1-01-1 p
1. --- Met de \VUgras{zomer} is het aanzicht van het land nog eens veranderd.
Het zaad dat aan de voor toevertrouwd was is gekiemd; een groen plantje kwam uit
de grond en schoot in de aar. Het graan werd gevormd, daarna rijpte het en
voortaan hoeft men alleen nog te oogsten.

%%K t08-c1-02-1 p
2. --- De \VUgras{veldbloemen} zijn opengegaan. \VUgras{Korenbloemen}
\textsuperscript{(1)}, \VUgras{klaprozen} \textsuperscript{(2)} en
\VUgras{vingerhoedskruid} \textsuperscript{(3)} mengen in de eentonige tint van
de tarwevelden de vrolijke tegenstelling van hun frisse \VUgras{bloemkronen}.

%%K t08-c1-03-1 p
3. --- De vruchtbomen hebben zich met bladeren getooid; de vrucht is gerijpt. De
kleine \VUgras{kersenboom} \textsuperscript{(4)} hangt vol mooie
\VUgras{kersen} \textsuperscript{(5)} die de kinderen met veel plezier plukken
en eten.

%%K t08-c1-04-1 p
4. --- Voortaan kan het vee de hele dag in de open lucht doorbrengen.

%%K t08-c1-04-2 p
De \VUgras{lammeren} \textsuperscript{(6)} zijn gegroeid en zijn mooie
\VUgras{ooien} \textsuperscript{(7)} en mooie \VUgras{rammen}
\textsuperscript{(8)} met een dichte vacht geworden. Zij grazen rustig het
\VUgras{gras} van de \VUgras{weide} \textsuperscript{(9)} die met
\VUgras{madeliefjes} gespikkeld is.

%%K t08-c1-04-3 p
Weldra zal men die dieren gaan scheren om hun \VUgras{wol} te krijgen waarvan
men het laken van onze kleren maakt.

%%K t08-tit-2 sub
\VUsaut{5.0mm}
\VUpk{10.2pt}{40}{De herder.}
\VUsaut{3.4mm}

%%K t08-c1-05-1 p
5. --- De \VUgras{herder} \textsuperscript{(10)} schijnt niet erg op hen te
letten. Hij bekommert zich alleen om het onweer dat aan de horizon voorbijtrekt,
en de \VUgras{schaapherder} \textsuperscript{(13)}, die een \VUgras{veldfles}
\textsuperscript{(14)} aan zijn lippen zet lijkt nog zorgelozer.

%%K t08-c1-06-1 p
6. --- De hitte is drukkend; want de \VUgras{hond} \textsuperscript{(15)} komt,
ook hij, zijn dorst lessen; hij likt het frisse water van het \VUgras{beekje}
\textsuperscript{(16)} en jaagt een arm \VUgras{kikkertje}
\textsuperscript{(17)} schrik aan dat zich in de oevergrassen verscholen had. Dit
springt in het heldere water waar de \VUgras{waterlelies}
\textsuperscript{(18)} bloeien.

%%K t08-c1-07-1 p
7. --- Een \VUgras{kwikstaart} \textsuperscript{(19)} baadt zich bij het
\VUgras{bronnetje} \textsuperscript{(20)} dat tussen twee rotsen opspringt en
verbergt zich onder de \VUgras{doornstruiken} \textsuperscript{(21)} en de
\VUgras{meidoornstruiken} \textsuperscript{(22)}.

%%K t08-tit-3 sub
\VUsaut{5.0mm}
\VUpk{10.2pt}{40}{De oogst.}
\VUsaut{3.4mm}

%%K t08-c1-08-1 p
8. --- Voor de boerenkinderen is de graanoogst een zeer vermakelijke tijd. Zij
maken vrolijk \VUgras{buitelingen} \textsuperscript{(24)} op de
\VUgras{strohopen} \textsuperscript{(23)}, zij helpen de \VUgras{oogsters}
\textsuperscript{(26)} en terwijl dezen het \VUgras{graanveld}
\textsuperscript{(27)} maaien met hun \VUgras{zeisen} \textsuperscript{(25)} die
met de \VUgras{wetsteen} \textsuperscript{(30)} goed gescherpt zijn, verzamelen
zij het graan en binden het tot \VUgras{schoven} \textsuperscript{(31)} of tot
\VUgras{bundeltjes} \textsuperscript{(32)}.

%%K t08-c1-09-1 p
9. --- Soms staat men de grootsten onder hen toe met een \VUgras{sikkel}
\textsuperscript{(12)} de \VUgras{gouden halmen} \textsuperscript{(28)} af te
snijden die de zware \VUgras{aren} \textsuperscript{(29)} vol korrels dragen; en
de kleinen, die op het \VUgras{vee} \textsuperscript{(33)} passen dat in de
\VUgras{weide} \textsuperscript{(34)} graast onder de schaduw van de grote
\VUgras{linde} \textsuperscript{(35)} vangen met hun \VUgras{net}
\textsuperscript{(36)} de mooie \VUgras{vlinders}.

%%K t08-c1-09-2 p
Sommigen klimmen zelfs op de \VUgras{kar} \textsuperscript{(37)} om de schoven
te schikken die men hun met een \VUgras{vork} \textsuperscript{(38)} aanreikt.

%%K t08-c1-10-1 p
10. --- Op de hele \VUgras{vlakte} \textsuperscript{(39)} waar de
\VUgras{leeuweriken} \textsuperscript{(41)} beginnen op te vliegen heerst de
bedrijvigheid. Allen zijn met de oogst bezig. Maar terwijl de armen doorgaan met
het oude gereedschap te gebruiken, \VUgras{zeisen, sikkels} en
\VUgras{dorsvlegels} \textsuperscript{(40)}, laten de rijke grondbezitters hun
tarwe of hun \VUgras{rogge} \textsuperscript{(43)} maaien met een
\VUgras{maaimachine} \textsuperscript{(42)}. Daarna maakt men, omdat men de
schoven niet naar de schuur hoeft te brengen, grote \VUgras{schovenhopen}
\textsuperscript{(44)}.

%%K t08-c1-11-1 p
11. --- Dan brengt men een \VUgras{dorsmachine} \textsuperscript{(47)} aan die
een \VUgras{locomobiel} \textsuperscript{(45)} met een \VUgras{drijfriem}
\textsuperscript{(46)} aandrijft en zo wordt het graan van het \VUgras{stro}
gescheiden, zonder het veld te verlaten waar het gezaaid werd.

%%K t08-tit-4 sub
\VUsaut{5.0mm}
\VUpk{10.2pt}{40}{Het onweer.}
\VUsaut{3.4mm}

%%K t08-c1-12-1 p
12. --- Dikwijls komen er zware \VUgras{onweders} \textsuperscript{(53)} in de
oogsttijd. Eerst hoort men van verre het gerommel van de \VUgras{donder}; een
\VUgras{hevige westenwind} \textsuperscript{(54)} begint te blazen en buigt
hijgend de dikste bomen van het \VUgras{bos} \textsuperscript{(49)}. Zwarte
\VUgras{wolken} \textsuperscript{(56)} bestormen de hemel; zij worden
doorkruist door verblindende zigzaggen van de \VUgras{bliksem}
\textsuperscript{(55)}. De \VUgras{regen} en soms de \VUgras{hagel} beginnen te
vallen, en slaan de oogst plat die klaar is om gemaaid te worden.

%%K t08-c1-13-1 p
13. --- Dikwijls steekt de bliksem een gebouw in brand. Zo werd vorig jaar het
dak van de \VUgras{windmolen} \textsuperscript{(50)} in de as gelegd en twee van
zijn \VUgras{wieken} \textsuperscript{(51)} vernield.

%%K t08-c1-14-1 p
14. --- Na enkele uren keert de rust terug. Een schitterende \VUgras{regenboog}
\textsuperscript{(52)} verschijnt aan de horizon en de natuur hervat haar leven
dat enkele ogenblikken onderbroken was. De landlieden, die in de \VUgras{hutten}
\textsuperscript{(48)} gevlucht waren gaan weer aan het werk en trachten moedig
de schade te herstellen die zij zojuist geleden hebben.

%%K t08-tit-5 sub
\VUsaut{6.0mm}
\VUcentre{13.2pt}{90}{\textit{Tweede tafereel.}}
\VUsaut{3.6mm}

%%K t08-tit-6 sub
\VUpk{10.2pt}{40}{De jacht. --- De wijnoogst.}
\VUsaut{1.4mm}
\VUpk{10.2pt}{40}{Het vissen.}
\VUsaut{4.0mm}

%%K t08-c2-01-1 p
1. --- Omstreeks het einde van \VUgras{augustus} of het midden van
\VUgras{september}, naar gelang van de streken, vindt de opening van de jacht
plaats. Nauwelijks hebben de eerste zonnestralen de \VUgras{hagedissen}
\textsuperscript{(2)} uit hun hol laten komen, of de \VUgras{jager}
\textsuperscript{(5)} gaat op weg met zijn \VUgras{honden}
\textsuperscript{(4)}.

%%K t08-c2-02-1 p
2. --- Hij heeft zijn stevige \VUgras{jachtkostuum} \textsuperscript{(9)} van
dik laken aangetrokken, leren \VUgras{beenkappen} aangedaan, waarmee hij erin
zal slagen te lopen in het vochtige gras waar de \VUgras{padden}
\textsuperscript{(1)} zich verbergen; hij heeft zijn \VUgras{geweer}
\textsuperscript{(6)} en zijn \VUgras{weitas} \textsuperscript{(10)} genomen, en
ziehier dat hij vertrekt.

%%K t08-c2-03-1 p
3. --- De ijver van de achtervolging brengt hem midden in een wijngaard waar de
wijnoogst zeer levendig is. De kinderen van de wijnbouwer, die gewoonlijk op de
geiten langs de weg passen, laten ze vandaag op goed geluk grazen, en nadat zij
een oude \VUgras{bok} \textsuperscript{(3)} aan een paal gebonden hebben die niet
zou nalaten van zijn vrijheid gebruik te maken om te ontsnappen, kennen zij zich
op hun beurt een aandeel in het genoegen toe. De een grijpt een druiventros in
een \VUgras{oogstmand} \textsuperscript{(12)} en vermaakt zich met het
\VUgras{sap} \textsuperscript{(14)} in een bekertje \textsuperscript{(13)} te
laten lopen om most (niet gegiste wijn) te maken. Een ander tooit zich met een
\VUgras{bladerenkrans} \textsuperscript{(15)} die hij zojuist gevlochten heeft.
Naast hen komen brutale \VUgras{mussen} \textsuperscript{(16)} zelfs voor de pers
de druiven roven.

%%K t08-c2-04-1 p
4. --- Terwijl de kinderen zich vermaken, werken de mannen. De
\VUgras{wijnplukkers} \textsuperscript{(18)} brengen de druiven in de kelder met
\VUgras{rugmanden} \textsuperscript{(17)}; een \VUgras{kuiper}
\textsuperscript{(19)} herstelt de \VUgras{vaten} \textsuperscript{(20)},
controleert of de \VUgras{duigen} \textsuperscript{(22)} en de
\VUgras{bodems} \textsuperscript{(21)} in goede staat zijn, en vervangt de
gebroken \VUgras{hoepels} \textsuperscript{(23)}. Hij heeft ook aan de
\VUgras{kuipen} \textsuperscript{(25)} gewerkt waarin men de \VUgras{druiven}
\textsuperscript{(26)} giet om ze te laten gisten.

%%K t08-c2-05-1 p
5. --- Wanneer de gisting geëindigd is, tapt men de wijn af en men legt het
overblijfsel op de \VUgras{pers} \textsuperscript{(28)} en door geleidelijk met
de \VUgras{grote schroef} \textsuperscript{(29)} aan te drukken perst men het
sap uit dat in een \VUgras{kuipje} \textsuperscript{(32)} loopt en men krijgt
nog \VUgras{wijn} \textsuperscript{(31)}.

%%K t08-tit-8 sub
\VUsaut{6.0mm}
\VUpk{10.2pt}{40}{Bier en cider.}
\VUsaut{3.4mm}

%%K t08-c2-06-1 p
6. --- Tegen de muur van de \VUgras{kelder} \textsuperscript{(27)} klimt een tak
\VUgras{hop} \textsuperscript{(30)}. Daar is hij alleen een sierplant, maar in
sommige landen, waar de wijnstok zeer zeldzaam is, wordt de hop geteeld om
\VUgras{bier} te maken.

%%K t08-c2-07-1 p
7. --- De kleine \VUgras{appelboom} \textsuperscript{(33)} is beladen met appels
die, wanneer zij rijp zullen zijn, een uitstekende \VUgras{cider} zullen
opleveren. In hun \VUgras{kooi} \textsuperscript{(34)} kijken de
\VUgras{kanarie} \textsuperscript{(35)} en de \VUgras{putter}
\textsuperscript{(36)} met jaloerse ogen naar de mussen die vrij dartelen.

%%K t08-c2-08-1 p
8. --- Zover het oog reikt, staan op de helling van de heuvels
\VUgras{wijnstokken} \textsuperscript{(40)} op een rij die de prachtige
\VUgras{stammen} \textsuperscript{(39)} dragen, getooid met zware
\VUgras{trossen}.

%%K t08-c2-08-2 p
Het hele jaar door heeft de \VUgras{wijnbouwer} \textsuperscript{(42)} gewerkt
om zijn \VUgras{wijngaard} \textsuperscript{(38)} te verzorgen en nu wordt hij
voor zijn moeite beloond met een mooie oogst. De \VUgras{wijnpluksters}
\textsuperscript{(41)} vullen de kuipen die op de kar staan die door
\VUgras{muildieren} \textsuperscript{(43)} getrokken wordt, en allen zijn vrolijk.

%%K t08-tit-9 sub
\VUsaut{5.0mm}
\VUpk{10.2pt}{40}{Het vissen.}
\VUsaut{3.4mm}

%%K t08-c2-09-1 p
9. --- Zo is het ook met de \VUgras{hengelaars} \textsuperscript{(44)} die zich
sinds de ochtend aan de oever van het water met hun \VUgras{hengels}
\textsuperscript{(45)} kwamen installeren.

%%K t08-c2-09-2 p
De \VUgras{vissen} \textsuperscript{(47)} bijten in de \VUgras{haak} zodra zij
hun \VUgras{lijn} \textsuperscript{(46)} in het water dompelen, en nauwelijks
hadden zij tijd om het \VUgras{aas} te vernieuwen, of een nieuw slachtoffer
spartelt aan het uiteinde van de lijn.

%%K t08-c2-09-3 p
Toch vangt de \VUgras{netvisser} \textsuperscript{(48)} die, vanuit zijn
\VUgras{boot} \textsuperscript{(49)}, het \VUgras{werpnet} midden in de
\VUgras{rivier} \textsuperscript{(50)} werpt meer vissen.

%%K t08-c2-10-1 p
10. --- De herfst is het seizoen van de mooie wandelingen : te voet, te
\VUgras{paard} \textsuperscript{(52)}, op de \VUgras{fiets}
\textsuperscript{(53)}, met de \VUgras{auto} \textsuperscript{(54)}. De
\VUgras{wegen} \textsuperscript{(51)} zijn schoon en men geniet van de laatste
mooie dagen, want ziehier dat de \VUgras{zwaluwen} \textsuperscript{(56)} zich
al op de daken verzamelen om met alle vleugels naar warmere landen te vliegen.
Het loof van de oude \VUgras{notenboom} \textsuperscript{(57)} wordt geel, de
\VUgras{noten} vallen en de hoge \VUgras{bomen} die als een \VUgras{bosje}
\textsuperscript{(58)} gegroepeerd staan op de \VUgras{hoogte}
\textsuperscript{(59)} zullen weldra hun bladeren verliezen.

%%K t08-c2-11-1 p
11. --- De \VUgras{zon} \textsuperscript{(60)} komt later op, de dagen worden
korter, een vrij scherpe kilte begint gevoeld te worden, in de \VUgras{ochtend},
en ziehier dat al vluchten \VUgras{houtsnippen} \textsuperscript{(61)} onze
ongastvrije streken verlaten.
\VUsaut{6.0mm}
\VUfilet{22mm}
